\section{Physical domain and topology transfer under containment}
\label{sec:physical-topology}

This section separates three claims that must not be conflated: pointwise
rank at a true cut, generic detection of a fixed noncut, and equality of cut
sets under strong-class containment.

\begin{lemma}[Strict isotropic subdivision]\label{lem:subdivision}
For \(y=(c,g,t)\in\Dplus\), put

\[
 B_+(y)=\max\{c,g,t,g+t-c,c+t-g,c+g-t\}<1.
\]

For every \(r\in(B_+(y),1)\),

\[
 (c,g,t)=(r,r,r)\odot(c/r,g/r,t/r)
\]

and both factors lie in \(\Dplus\).  In strict continuous time the same holds
whenever

\[
 r> B_{\mathrm{CT}}(y)
 :=\max\{c,g,t,gt/c,ct/g,cg/t\}.
\]
\end{lemma}

\begin{proof}
The first six lower bounds are precisely the upper-coordinate and three
inverse-Fourier inequalities for the residual triple.  The isotropic factor
is strict for \(0<r<1\).  In continuous time, for example,
\(c/r>(g/r)(t/r)\) is equivalent to \(r>gt/c\); the other inequalities are
cyclic.  All bounds are strictly below one at an interior point.
\end{proof}

\begin{lemma}[Root movement]\label{lem:root-movement}
Moving the root to any admissible rooting of the same standard semi-directed
mixed graph leaves the K3P physical image and its regular germs unchanged up
to analytic physical reparameterization.
\end{lemma}

\begin{proof}
K3P transition matrices are symmetric and the root distribution is uniform,
so reversing ordinary arcs on an old-to-new root path preserves every
displayed-tree probability.  Use \cref{lem:subdivision} at the new root and
suppress the old one.  If the old root meets a reticulation-parent edge,
inspect each switching: the selected parent uses only the serial product,
whereas the unselected one-child stem subtends all observed leaves and has
total character zero.  Every summand of \eqref{eq:k3p-fourier} is unchanged
after transporting the reversible edge products.  More explicitly, write the
two parent weights at a reticulation in the recorded parent order as
\((\lambda,1-\lambda)\).  Root movement preserves \(\lambda\) when that order
is preserved and transports it to \(\lambda'=1-\lambda\) when the order is
reversed; thus the same switching keeps the same coefficient.  Strict
subdivision supplies analytic physical
sections in both directions on a neighborhood, which is the asserted
reparameterization.
\end{proof}

For a split \(S=A\mid A^c\), order the Fourier flattening by the bridge
character \(h=\xor_{i\in A}h_i\).  It is block diagonal with four character
blocks.

\begin{proposition}[True-cut rank, pointwise]\label{prop:true-cut}
At every strict K3P parameter point of every network, if \(S\) is a bridge
split, then

\[
 \rank\Flat_S(q)\le4.
\]
\end{proposition}

\begin{proof}
Condition on the character crossing the bridge.  Each of the four character
blocks is a positive rank-one outer product of the two side tensors, with the
corresponding bridge multiplier in the nonzero blocks.  Summing their ranks
gives at most four.
\end{proof}

\begin{lemma}[Displayed-tree witness for a noncut]
\label{lem:displayed-tree-noncut}
If \(S=A\mid A^c\) is not a bridge split of a strongly tree-child level-2
network, some displayed switching gives a tree that does not display \(S\).
\end{lemma}

\begin{proof}
Color the leaves by the two sides of \(S\), and take the two color hulls in
the network's abstract bridge-incidence tree.  Disjoint hulls would be
separated by a bridge realizing \(S\), contrary to hypothesis.  If their
intersection contains an edge, that edge is a bridge split crossing \(S\).
Every displayed tree retains this bridge split, and two splits of one tree
are compatible, so no displayed tree can also display \(S\).

Otherwise the hulls meet in one complete component.  Every incident branch
there is monochromatic, and minimality of the two hulls supplies at least two
branches of each color.  The balanced compression argument of
\cref{lem:balanced-compression} applies inside this cycle or theta factor and
produces a switching with a wrong two-by-two quartet.  Extending its local
choices arbitrarily over all other reticulations gives the required global
displayed tree.
\end{proof}

\begin{proposition}[Generic noncut recovery]\label{prop:generic-noncut}
Fix a standard strongly tree-child level-2 network \(N\) and a split \(S\)
that is not a bridge split.  Some \(5\)-by-\(5\) minor of
\(\Flat_S\circ\Phi_N\)
is a nonzero polynomial.  Consequently
\(\rank\Flat_S(q)>4\) on a Zariski-open dense subset of the physical parameter
domain.
\end{proposition}

\begin{proof}
Choose the switching in \cref{lem:displayed-tree-noncut}.  The ordinary-tree
quartet criterion supplies \(a_0,a_1\in A\) and
\(a_2,a_3\in A^c\) whose restricted displayed tree has a quartet split
different from \(a_0a_1\mid a_2a_3\).  Work first at the polynomial-map
boundary by setting each inheritance probability to zero or one so that only
this switching remains, and give every edge an isotropic multiplier strictly
between zero and one.

Relabel the displayed quartet split as \(01\mid23\) and its wrong flattening
as \(02\mid13\).  After suppressing its serial paths, write
\(p_0,p_1,p_2,p_3\in(0,1)\) for the four pendant multipliers and
\(u\in(0,1)\) for the internal multiplier.  In the zero-character block, the
rows and columns indexed by \(00\) and \(CC\) give determinant

\[
 p_0p_1p_2p_3(1-u^2)>0.
\]

One positive entry from each of the other three character blocks augments
this to a nonzero \(5\)-by-\(5\) minor.  Assigning character zero to every
unselected leaf identifies it with a minor of \(\Flat_S\circ\Phi_N\).
Thus that network minor has a nonzero boundary evaluation and is not the zero
polynomial.  By continuity it remains nonzero when every inheritance
probability is moved sufficiently close to its selected endpoint but into
\((0,1)\).  The isotropic edge points are strict in both \(\Dplus\) and
\(\Dct\).  Hence a strict physical evaluation exists, and the zero set of the
minor has empty interior in the open K3P parameter domain.
\end{proof}

\begin{corollary}[The easy directed cut inclusion]\label{cor:target-source-cut}
If \(N\preceqplus N'\), then every target bridge split is a source bridge
split:

\[
 \operatorname{Cut}(N')\subseteq\operatorname{Cut}(N).
\]
\end{corollary}

\begin{proof}
For a target bridge, all \(5\)-minors vanish identically after composition
with the physical target section.  They therefore vanish on the source-open
parameter neighborhood.  If the split were a source noncut,
\cref{prop:generic-noncut} would make one pullback a nonzero polynomial,
which cannot vanish on an open set.
\end{proof}

The reverse inclusion needs genuinely K3P pointwise information, but only for
a finite one-active universe.

\begin{lemma}[Balanced noncut compression]\label{lem:balanced-compression}
Let a complete strongly tree-child cycle or theta factor carry a two-color
boundary assignment with at least two actual labels of each color and no
bridge realizing the color split.  There is a restriction which retains the
primitive core, one complete minimum strong repair, every path-sink child,
and at least two actual labels of each color.  All required unselected roles
remain as zero-character completion ports.  The restricted coloring remains
noncut, some displayed switching fails it, and two actual labels of each
color give a wrong quartet in the finite one-active universe of
\cref{lem:one-active}.
\end{lemma}

\begin{proof}
First retain a rigid support \(Q\): every path-sink child and one occupied port
on every segment of one complete minimum repair.  The core table in
\cref{sec:bounded} gives \(|Q|\le4\), and suppressing unselected ordinary
subdivisions leaves the primitive cycle or theta core.  A nonroot incoming
boundary, when present, is retained in addition; because it is already an
actual colored label, the number of further representatives needed to reach
two of each color decreases by one.  Thus the compulsory support plus the
four-active representatives still has at most eight ports.  On each directed
segment, read the colors as a word and decompose it into maximal constant runs.

If some segment contains three consecutive runs \(c,d,c\), retain one actual
label from each of those runs and, if necessary, one further \(d\)-label from
the same or another segment.  In every switching the path between the two
outer \(c\)-attachments contains the middle \(d\)-attachment.  No tree edge
can separate the complete color classes: a segment edge separates the outer
\(c\)-ports or retains the middle \(d\) with them, a pendant edge isolates one
port, and another \(d\)-port exists.  This is a direct all-switching
obstruction.

Otherwise every segment has at most two runs.  Select actual labels from at
most two runs of each color, preferring labels already in \(Q\), until two
labels of each color are retained.  If a color is represented by only one
selected label, retain a second actual label in the same run when possible,
and otherwise from another run or segment.  Its primitive blob and all
compulsory child roles remain; no internal blob edge is a bridge, and a
pendant arm cannot separate a color class containing two labels.  Thus the
restricted split is still noncut.

It remains to pass from this bounded restriction to the exact four-active
palette.  Retain one actual label from each selected run, preserving every
occupied segment and fixed child role.  Each segment word is now one of
\[
 (),\quad(0),\quad(1),\quad(0,1),\quad(1,0).
\]
Double a singleton run when needed to keep two actual labels of its color.
Any split displayed by every switching before this reduction would remain so
after leaf restriction.  The independent finite switching replay has zero
survivors in this direct-or-singleton-doubled palette, so in this second case
some switching fails the color split.  In either case choose a failing
switching.  The ordinary-tree quartet criterion then selects two retained
actual labels of each color whose quartet fails that split.  Every
other repair or path-sink role is retained at character zero, so the result
is one of the graph-derived one-active directions rather than a pruned
topology.

The exact censuses are reproducibility checksums, not premises.  A standalone
word program exhausts all (808{,}642) balanced distributions with four
through eight active ports and partitions them into (229{,}988) three-run
obstructions, (544{,}350) direct-palette reductions, and (34{,}304)
singleton-doubled reductions.  A no-import implementation independently
constructs the five primitive cores and tests (379{,}742) valid compressed
presentations, with zero all-switching survivors.  Finally, five cores give
72 active-labelled records and 216 split entries; 12 are displayed by every
switching, leaving the 204 wrong directions certified in
\cref{lem:one-active}.
\end{proof}

\begin{lemma}[One-active K3P obstruction]\label{lem:one-active}
For a balanced two-color split across one active strongly tree-child
level-2 component, primitive-core and word compression produce exactly 216
labelled split entries.  Exactly 12 are displayed by every switching.
For each of the remaining 204 normalized wrong-split directions, every
strict \(\Dplus\) tensor has flattening rank greater than four.
\end{lemma}

\begin{proof}[Exact finite proof]
The active model-independent graph producer derives five primitive core
templates, 72 active-labelled records, and three \(2\)-by-\(2\) splits per
record from literal directed graphs, root suppression, switching choices, and
descendant masks.  Its fresh certificate is required to agree byte for byte
with the exact topology input consumed by the K3P replay.  Starting from those
independently regenerated switching signatures, that replay recomputes the
displayed-by-all flags: it removes 12 entries and gives 204 distinct labelled
direction keys, without quotienting by an automorphism.  The five-core theorem
and the compression proof are handwritten above; the 72-row labelled table is
also graph-derived active evidence, not a K3P algebraic conclusion.

For 180 directions one selected block minor has a strictly signed
Bernstein expansion on the strict K3P simplex.  Twelve more are excluded by
positive signed combinations of two minors.  Ten use cyclic six-minor
identities.  The last two use two independently checked cyclic elimination
certificates, one obtained by an exact labelled transport.  The disjoint
partition is

\[
 204=180+12+10+1+1.
\]

If the total flattening had rank at most four, its four positive character
blocks would each have rank one, forcing all selected minor equations to
vanish.  The corresponding signed or cyclic certificate contradicts that
system throughout \(\Dplus\).  A separately implemented verifier rebuilds all
204 normalized directions and all exact K3P arithmetic from that table.  The
graph producer, byte comparison, and coherent topology/mask mutations are
separate active gates.  The supplement records this implementation boundary.
\end{proof}

\begin{theorem}[Strong-class cut-set equality under containment]
\label{thm:cut-transfer}
If \(N\preceqplus N'\) for networks in the class of \cref{thm:main}, then

\[
 \operatorname{Cut}(N)=\operatorname{Cut}(N').
\]

Hence the labelled abstract bridge-incidence trees agree.  This conclusion
does not yet decide whether a trivalent internal component is ordinary or is
the contraction of a three-boundary cycle; that decoration is recovered in
\cref{lem:trivalent-decoration}.
\end{theorem}

\begin{proof}
By \cref{cor:target-source-cut}, only
\(\operatorname{Cut}(N)\subseteq\operatorname{Cut}(N')\) remains.  Suppose a
nontrivial source bridge split \(S=A\mid B\) is absent from the target; a
singleton split is pendant in every standard network and cannot be lost, so
\(|A|,|B|\ge2\).
Color target leaves by \(A,B\), and in the target reduced tree let \(H_A,H_B\)
be the two color hulls.  They intersect, because an edge separating disjoint
hulls would realize \(S\) as a target bridge.

If \(H_A\cap H_B\) contains an edge, its target bridge split \(R=C\mid D\)
has a leaf in each of
\(A\cap C,A\cap D,B\cap C,B\cap D\); thus \(R\) crosses \(S\).  But
\cref{cor:target-source-cut} makes \(R\) a source bridge.  Two edge splits of
one tree are compatible, so source bridge splits \(R,S\) cannot cross.  This
excludes the two-active alternative without assuming a common bridge tree.

Otherwise \(H_A\cap H_B=\{v\}\) is one target component.  Every branch of
the target reduced tree at \(v\) is monochromatic.  Indeed, if one branch
contained both colors, then, because \(v\) lies in each color hull, the edge
entering that branch would belong to both \(H_A\) and \(H_B\), contrary to
their intersection being only \(v\).  Moreover, membership of \(v\) in each
minimal color hull forces at least two incident branches of each color: a
labelled tree leaf has only one color and cannot lie in both hulls, so \(v\)
is internal in each hull.
Thus \(v\) has at least four relevant branches.  A binary ordinary tree
vertex has degree three, so \(v\) represents a nontrivial central blob.

Apply \cref{lem:balanced-compression} to this central blob.  It supplies
actual labels \(a_1,a_2\in A\) and \(b_1,b_2\in B\) whose induced quartet
fails to display \(a_1a_2\mid b_1b_2\), while the completion roles are
retained at zero character.  This is one of the 204 graph-derived labelled
wrong-split directions in \cref{lem:one-active}.

It remains to check that suppressing the material between the four selected
labels and the central blob stays inside the strict K3P domain.  A nonempty
ordinary serial path is a convolution of strictly positive
\(\mathbb Z_2^2\) transition distributions, hence is a strict K3P edge.
Every component strictly between \(v\) and one selected label has only two
retained boundaries.  Marginalizing such a noncentral blob gives

\[
 M_{\mathrm{eff}}=\sum_s w_sM_s,
 \qquad w_s>0,\qquad \sum_s w_s=1,
\]

where \(M_s\) is the strict K3P path matrix of switching \(s\).  Each \(w_s\)
is a product of local factors \(\lambda_r\) or \(1-\lambda_r\), so positivity
and normalization follow from summing all local switches.  K3P matrices form
a linear family, so \(M_{\mathrm{eff}}\) is K3P; its transition entries are
positive convex combinations of positive entries and its three nontrivial
Fourier eigenvalues are convex combinations of numbers in \((0,1)\).  It is
therefore a strict \(\Dplus\) edge.  Chains of these two-boundary mixtures
and ordinary paths remain strict by convolution.
Retained inheritance probabilities are \(\lambda\) or \(1-\lambda\), still
in \((0,1)\).  Thus the compressed target quartet lies at a strict physical
point of the 204-direction theorem; no target regularity or target-open
marginal is being assumed.

Apply the same four-leaf marginal to the containment identity.  In Fourier
coordinates this sets every omitted leaf character to zero.  On the source
side the original bridge \(S\) survives because two selected labels lie on
each side, so \cref{prop:true-cut} gives rank at most four.  On the identical
target tensor, \cref{lem:one-active} gives rank greater than four at every
strict target point.  This contradiction proves the reverse inclusion.  The
compatible split system determines the labelled abstract bridge-incidence
tree.  As usual for edge-split recovery, it does not by itself decorate a
trivalent vertex as an ordinary branching component or a three-boundary blob.
\end{proof}

\begin{remark}[Exact claim boundary]\label{rem:no-universal-cut}
Theorem \ref{thm:cut-transfer} is a directional theorem for two networks in
the strong class.  It does not promote \cref{prop:generic-noncut} to a
universal pointwise converse for arbitrary K3P networks.  No later proof uses
such a converse.
\end{remark}
